\documentclass[preview]{standalone}

\usepackage{amsmath}
\usepackage{amssymb}
\usepackage{stellar}
\usepackage{definitions}

\begin{document}

\id{psicologia-attenzione}
\genpage

\section{Attenzione}

\begin{snippetexample}{attenzione-festa-example}{Attenzione in una festa}
    Durante una festa rumorosa, una persona cerca un'amica in mezzo alla folla:
    individua il vestito rosa che indossa, sente la sua voce tra musica e
    conversazioni, si volta quando sente un rumore di vetri rotti e poi riprende
    la ricerca. Individuare l'amica, riconoscere una voce rilevante, orientarsi
    verso un rumore improvviso e tornare al compito precedente sono tutte
    operazioni che coinvolgono l'attenzione.
\end{snippetexample}

\begin{snippetdefinition}{attenzione-definition}{Attenzione}
    Dal punto di vista cognitivo, l'\emph{attenzione} è l'insieme dei meccanismi
    che permettono di concentrare e focalizzare le risorse mentali su alcune
    informazioni, definendo ciò di cui il soggetto è consapevole in un dato
    momento.
\end{snippetdefinition}

\begin{snippet}{valorizzazione-inibizione-attenzione}
    L'attenzione opera attraverso due funzioni complementari. La
    \emph{valorizzazione} seleziona alcune informazioni per sottoporle a
    elaborazioni ulteriori; l'\emph{inibizione} riduce invece l'elaborazione di
    informazioni considerate irrilevanti.
\end{snippet}

\begin{snippetdefinition}{attenzione-selettiva-definition}{Attenzione selettiva}
    L'\emph{attenzione selettiva} è il processo che seleziona alcune
    informazioni sulla base della loro salienza e le elabora con particolare
    efficienza.
\end{snippetdefinition}

\begin{snippet}{forme-attenzione-selettiva}
    L'attenzione selettiva può riferirsi a una modalità sensoriale, a una zona
    dello spazio, a uno stimolo specifico, a una caratteristica dello stimolo
    come colore o forma, oppure a un movimento. Per questo l'attenzione non è
    solo spaziale: spesso procede selezionando oggetti e proprietà rilevanti.
\end{snippet}

\includesnpt[width=62\%,src=/snippet/static-psychology/selective-attention-example.jpg]{centered-img}

\begin{snippetdefinition}{attenzione-basata-oggetti-definition}{Attenzione basata sugli oggetti}
    L'\emph{attenzione basata sugli oggetti} è una forma di attenzione per cui,
    quando l'attenzione è diretta verso un oggetto, tutte le parti di quell'oggetto
    vengono selezionate simultaneamente per essere elaborate.
\end{snippetdefinition}

\begin{snippetexample}{oggetto-attenzione-example}{Oggetto e attenzione}
    Se una persona si focalizza su un'amica vestita di rosa, ha una probabilità
    maggiore di notare l'orologio al suo polso, parte dello stesso oggetto
    percettivo, rispetto all'orologio indossato da una persona accanto, anche se
    i due orologi si trovano alla stessa distanza.
\end{snippetexample}

\begin{snippetdefinition}{fuoco-attenzione-definition}{Fuoco dell'attenzione}
    Il \emph{fuoco dell'attenzione} è la porzione di stimolazione o informazione
    su cui vengono concentrate le risorse attentive in un dato momento.
\end{snippetdefinition}

\begin{snippet}{ampiezza-fuoco-attenzione}
    Il fuoco dell'attenzione può essere ristretto o esteso. Quando è ristretto,
    l'elaborazione delle informazioni selezionate tende a essere più efficiente;
    quando è ampio e disperso, l'elaborazione tende a essere meno precisa.
\end{snippet}

\begin{snippetdefinition}{effetto-cocktail-party-definition}{Effetto cocktail party}
    L'\emph{effetto cocktail party} è il fenomeno per cui una persona riesce a
    percepire informazioni rilevanti, come il proprio nome, in un ambiente
    rumoroso, anche quando è impegnata in un'altra conversazione.
\end{snippetdefinition}

\begin{snippet}{broadbent-filtro-attentivo}
    Secondo Donald Broadbent, la mente opera solo su alcuni stimoli in entrata
    grazie a un sistema di filtraggio. L'effetto cocktail party limita però
    l'idea di un filtro completamente precoce, perché mostra che informazioni
    rilevanti possono essere rilevate anche quando non sono al centro
    dell'attenzione.
\end{snippet}

\begin{snippetdefinition}{attenzione-divisa-definition}{Attenzione divisa}
    L'\emph{attenzione divisa} è il processo attraverso cui l'attenzione viene
    distribuita tra diversi stimoli o compiti a partire da una quantità limitata
    di risorse cognitive.
\end{snippetdefinition}

\begin{snippetexample}{esperimento-duncan-example}{Esperimento di Duncan}
    In un esperimento di Duncan, i partecipanti vedevano al centro di un monitor
    un rettangolo con una piccola apertura su un lato e una linea che lo
    attraversava. Quando dovevano rispondere a due domande relative allo stesso
    oggetto, per esempio se il rettangolo fosse grande o piccolo e se
    l'apertura fosse a sinistra o a destra, l'accuratezza era elevata. Quando
    invece dovevano rispondere a domande relative sia al rettangolo sia alla
    linea, l'accuratezza diminuiva. Il risultato suggerisce che l'attenzione è
    più efficiente quando le informazioni appartengono allo stesso oggetto.
\end{snippetexample}

\includesnpt[width=55\%,src=/snippet/static-psychology/duncan-divided-attention-stimuli.jpg]{centered-img}

\begin{snippet}{competizione-attentiva}
    Secondo Duncan, l'attenzione è l'esito di una competizione tra stimoli che
    può verificarsi in diverse fasi dell'elaborazione. Lo stimolo che riceve
    più risorse viene analizzato in modo più completo rispetto agli altri.
\end{snippet}

\end{document}
